\documentclass[language=ngerman,algolist]{TemplateLecture}

\renewcommand{\LectureName}{Randomisierte und Approximative Algorithmen}
\renewcommand{\ProfName}{Prof. Dr. Thomas Kesselheim, \\ Prof. Dr. Heiko Röglin}
\renewcommand{\ProfNameShort}{Profs. Kesselheim, Röglin}
\renewcommand{\Semester}{Winter 2025/26}
\renewcommand{\mName}{Daniel Emse}


\begin{document}
\setcounter{chapter}{0}

%%%%%%%%%%%%%%%%%%%%%%%%%%%%%%%%%%%%%%%%%%%%%%%%%%%%%%%%%%%%%%%%%%%%%%%%%%%%%%%%%%%%%%%%%%%%%%%%%%%%
%%%%%%%%%%%%%%%%%%%%%%%%%%%%%%%%%%%%%%%%%% 2025 / 10 / 14 %%%%%%%%%%%%%%%%%%%%%%%%%%%%%%%%%%%%%%%%%%
%%%%%%%%%%%%%%%%%%%%%%%%%%%%%%%%%%%%%%%%%%%%%%%%%%%%%%%%%%%%%%%%%%%%%%%%%%%%%%%%%%%%%%%%%%%%%%%%%%%%


\part{Randomisierte und Approximative Algorithmen}

\chapter{Einleitung}
\begin{defi}{Optimierungsproblem, Approximationsalgorithmus}{optimizationProblemApproxAlgo}
    Ein \emph{Optimierungsproblem} $\Pi$ besteht aus
    \begin{itemize}
        \item einer Meinge von Instanzen $\cI_\Pi$
        \item für jede Instanz $ I \in \cI_\Pi$
        \begin{itemize}
            \item einer Menge $\cS_I$ von Lösungen
            \item einer Zielfunktion $f_I: \cS_I \to \IR_{\geq 0}$.
        \end{itemize} 
    \end{itemize}
    Ein \emph{Approximationsalgorithmus} nimmt $I$ als Eingabe und berechnet $A(I) \in \cS_I$ (in polynomieller Zeit).

    Sei $w_A(I) = f_I(A(I))$. Ein Approximationsalgorithmus erreicht Approximationsfaktor $r$, falls 
    \begin{equation*}
        w_A(I) \leq r \cdot \OPT(I) \quad \forall I \in \cI_\Pi \quad \quad \text{(für Minimierungsprobleme)},
    \end{equation*}
    bzw.
    \begin{equation*}
        w_A(I) \geq \frac{1}{r} \cdot \OPT(I) \quad \forall I \in \cI_\Pi \quad \quad \text{(für Maximierungsprobleme)}.
    \end{equation*}
\end{defi}

\chapter{Greedy-Algorithmen}

\section{Vertex Cover}
\newLecture{14.10.2025}


\begin{prob}{Vertex Cover}{VertexCover}
    \prblmObj{Ungerichteter Graph $G = (V, E)$}{Alle $V' \subseteq V$, sodass $\forall e = (u, v) \in E$ gilt $u \in V'$ oder $v \in V'$}{Minimiere $\abs{V'}$}\index{Vertex Cover}
\end{prob}




\begin{defi}{Matching}{Matching}
    Ein \emph{Matching}\index{Matching} $M \subseteq E$ eines Graphen $G = (V, E)$ ist eine Teilmenge der Kanten, sodass es keinen Knoten gibt, der zu mehr als einer Kante aus $M$ inzident ist.

    Es heißt \emph{inklusionsmaximal}\index{Matching!inklusionsmaximal}, wenn für alle Kanten $e \in E \setminus M$ gilt: $M \cup \set{e}$ ist kein Matching.
\end{defi}
\begin{algorithm}[h]
    \caption{\textsc{Matching-VC}}\label{alg:matching-vc}\index{Matching-VC}
    \begin{algorithmic}[1] % Matching-VC
        \Input{ungerichteter Graph $G = (V, E)$}
        \Output{Vertex Cover $V' \subseteq V$ mit $\abs{V'} \leq 2 \cdot \OPT(G)$}
        \LComment{Berechne inklusionsmaximales Matching $M \subseteq E$:}
        \State{} $M := \emptyset$
        \ForAll{$e = \set{u, v} \in E$}
            \If{weder $u$ noch $v$ ist inzident zu einer Kante in $M$}
                \State{} $M := M \cup \set{e}$
            \EndIf{}
        \EndFor{}
        \State{} $V(M) := \set{v \in V \mid \text{$v$ ist inzident zu einer Kante in $M$}}$
        \State{} \Return{$V(M)$}
    \end{algorithmic}
\end{algorithm}


\begin{thm}{Matching-VC}{MatchingVC}
    Matching-VC hat Laufzeit $\cO(\abs{V} + \abs{E})$ für Graphen in Adjazenzlistendarstellung. Er berechnet ein Vertex Cover der Größe höchstens $2 \cdot \OPT(G)$, wobei $\OPT(G)$ die Größe des kleinsten Vertex Covers ist.
\end{thm}

\begin{proof}
    \begin{description}
        \item[Laufzeit:] In Schritt 1 wird jede Kante genau einmal betrachtet, $\cO(\abs{V} + \abs{E})$. 
        
        In Schritt 2 wird jeder Knoten höchstens einmal betrachtet, $\cO(\abs{V})$. Gemeinsam ergibt sich $\cO(\abs{V} + \abs{E})$.
        \item[Korrektheit:] Angenommen, $e = (u,v) \in E$ wurde nicht von $V(M)$ abgedeckt. Das heißt, es wären weder $u \in V(M)$ noch $v \in V(M)$. Dann wäre $M \cup \set{e}$ ein Matching und $e$ würde in Schritt 1 zu $M$ hinzugefügt werden. Widerspruch zur Inklusionsmaximalität von $M$. Also ist $V(M)$ ein Vertex Cover.
        \item[Approximation:] Sei $V^*$ ein minimales Vertex Cover. $V^*$ muss alle Kanten $e \in E$ abdecken, also auch alle $e \in M$. 
        
        Außerdem haben in $M$ keine zwei Kanten denselben Endpunkt. Deshalb ist $\abs{M} \leq \abs{V^*}$. 
        
        Für unsere Lösung gilt $\abs{V(M)} = 2 \cdot \abs{M} \leq 2 \cdot \OPT(G)$.
    \end{description}
\end{proof}

\section{Set Cover}

\begin{prob}{Set Cover}{SetCover}
    \prblmObj{Menge $S$ mit $n$ Elementen, $n$ Teilmengen $S_1, \ldots, S_m \subseteq S$ mit $\bigcup_{i=1}^m S_i = S$, Kosten $c_1, \ldots, c_m \in \IN$.}{Alle $A \subseteq \set{1, \ldots, m}$ mit $\bigcup_{i \in A} S_i = S$}{Minimiere $\sum_{i \in A} c_i$}\index{Set Cover}
\end{prob}


\begin{algorithm}[h] % Greedy-SC
    \caption{\textsc{Greedy-SC}}\label{alg:greedySC}\index{Greedy-SC}
    \begin{algorithmic}[1]
        \Input{Grundmenge $S$ mit $n$ Elementen, $n$ Teilmengen $S_1, \ldots, S_m \subseteq S$ mit $\bigcup_{i=1}^m S_i = S$, Kosten $c_1, \ldots, c_m \in \IN$.}
        \Output{$A \subseteq \set{1, \ldots, m}$ mit $\bigcup_{i \in A} S_i = S$ und $\sum_{i \in A} c_i$ minimal}
        \State{} $A := \emptyset, C := \emptyset$
        \While{$C \neq S$}
            \State{} Wähle $S_i$ mit minimalen relativen Kosten $r_i(C) = \frac{c_i}{\abs{S_i \setminus C}}$
            \State{} Setze $\price(x) := r_i(C)$ für alle $x \in S_i \setminus C$
            \State{} $A := A \cup \set{i}$
            \State{} $C := C \cup S_i$
        \EndWhile{}
        \State{} \Return{$A$}
    \end{algorithmic}
\end{algorithm}

\begin{thm}{Greedy-SC}{greedySC}
    \textsc{Greedy-SC} ist ein $H_n$-Approximationsalgorithmus für Set Cover, wobei
    \begin{equation*}
        H_n = \sum_{i = 1}^n \frac{1}{i}
    \end{equation*}
    die $n$-te harmonische Zahl ist.
\end{thm}

\emph{Untere Schranke für Greedy-SC:} Sei $\varepsilon > 0$ sehr klein. Betrachte $n$ Knoten und $m = n + 1$ Mengen, wobei $S_i = \set{i}$ mit Kosten $\frac{1}{n+1-i}$ ist (für $i \in S$) und $S_{n+1} = S$ mit Kosten $1 + \varepsilon$. Dann wählt \textsc{Greedy-SC} wählt alle Einzelnen, OPT die gesamte Menge.

\newLecture{15.10.2025}

\begin{proof}
    \emph{Beobachtung:} Es gilt 
    \begin{equation*}
        c(A) = \sum_{x \in S} \price(x).
    \end{equation*}
    Sei $x_1, \ldots, x_n$ die Reihenfolge der Elemente in $S$, wie sie vom \textsc{Greedy-SC} abgedeckt werden (Beim Abdecken von mehreren Elementen in gleicher Iteration wähle beliebige Reihenfolge).

    \begin{lem}{}{GreedySCPrice}
        Für alle $k \in \set{1,\ldots,n}$ gilt
        \begin{equation*}
            \price(x_k) \leq \frac{\OPT}{n-k+1}.
        \end{equation*}
    \end{lem}
    \begin{proof}[Beweis: Lemma~$\ref{lem:GreedySCPrice}$]
        Betrachte ein beliebiges $k \in \set{1,\ldots,n}$. Sei $i \in \set{1,\ldots,m}$ der Index der Menge $S_i$, mit der $x_k$ erstmalig abgedeckt wird. Bezeichne $A$ die Auswahl der Mengen, \emph{bevor} $x_k$ abgedeckt wird.

        Sei $C := \bigcup_{i \in A} S_i$. Es gilt $\abs{C} \leq k - 1$. Eine optimale Auswahl $A^*$ von Mengen deckt ganz $S$ ab, insbesondere auch $S \setminus C$, wobei $\abs{S \setminus C} \geq n - k + 1$.

        Es gilt 
        \begin{equation*}
            \begin{split}
                \OPT &= \sum_{j \in A^*} c_j \geq \sum_{\substack{j \in A^* \\ S_j \setminus C \neq \emptyset}} c_j = \sum_{\substack{j \in A^* \\ S_j \setminus C \neq \emptyset}} \abs{S_j \setminus C} \underset{=r_j(C) \geq r_i(C)}{\underbrace{\frac{c_j}{\abs{S_j \setminus C}}}}  \\
                &\geq r_i(C) \underset{\geq \abs{S \setminus C}}{\underbrace{\sum_{\substack{j \in A^* \\ S_j \setminus C \neq \emptyset}} \abs{S_j \setminus C}}} \geq \underset{=\price(x_k)}{\underbrace{r_i(C)}} \cdot (n-k+1)
            \end{split}
        \end{equation*}
        wobei $r_j(C) \geq r_i(C)$ wegen der Greedy-Auswahl und 
        \begin{equation*}
            \sum_{\substack{j \in A^* \\ S_j \setminus C \neq \emptyset}} \abs{S_j \setminus C} \geq \abs{S \setminus C},
        \end{equation*}
        weil $A^*$ ganz $S \setminus C$ abdeckt.
    \end{proof}
    Nun gilt mit Lemma~\ref{lem:GreedySCPrice}
    \begin{equation*}
        \begin{split}
            c(A) &= \sum_{k = 1}^n \price(x_k) \leq \sum_{k = 1}^n \frac{\OPT}{n-k+1} = \OPT\sum_{i=1}^n \frac{1}{i} = \OPT \cdot H_n
        \end{split}
    \end{equation*}
    und Satz~\ref{thm:greedySC} ist gezeigt.
\end{proof}

\section{Scheduling auf identischen Maschinen}
\begin{prob}{Scheduling auf identischen Maschinen}{SchedulingIdenticalMachines}\index{Scheduling}\index{Scheduling!auf identischen Maschinen}
    \prblmObj{Menge $J = \set{1,\ldots,n}$ von Jobs, Menge $M = \set{1,\ldots,m}$ von Maschinen, Größen $p_j > 0$ für jeden Job $j \in J$}{Schedule $\pi: J \to M$, Last von Maschine $i$ in Schedule $\pi$: $L_i(\pi) = \displaystyle\sum_{\substack{j \in J \\ \pi(j) = i}} p_j$, Makespan $C(\pi) = \max_{i \in M} L_i(\pi)$}{Finde Schedule $\pi$ mit kleinstem Makespan}
\end{prob}

\begin{algorithm}[H]
    \caption{\textsc{LeastLoaded}}\label{alg:leastloaded}\index{LeastLoaded}
    \begin{algorithmic}[1]
        \Input{Siehe~\Cref{prob:SchedulingIdenticalMachines}}
        \Output{Siehe~\Cref{prob:SchedulingIdenticalMachines}}
        \State{Betrachte Jobs in Reihenfolge $1,\ldots,n$}
        \State{Weise jedem Job der Maschine mit bislang kleinster Last zu}
    \end{algorithmic}
\end{algorithm}

\begin{thm}{LeastLoaded}{LeastLoaded}
    \textsc{LeastLoaded} ist ein $(2-\frac{1}{m})$-Approximationsalgorithmus für Scheduling auf identischen Maschinen.
\end{thm}

\begin{proof}
    Sei $\pi^*$ ein optimaler Schedule. Dann ist 
    \begin{equation*}
        C(\pi^*) \geq \frac{1}{m}\sum_{j \in J}  p_j,
    \end{equation*}
    denn
    \begin{equation*}
        C(\pi^*) = \max_{i \in M} \sum_{\substack{j \in J \\ \pi(j) = i}} p_j \geq \frac{1}{m}\sum_{i \in M}\sum_{\substack{j \in J \\ \pi(j) = i}} = \frac{1}{m}\sum_{j \in J} p_j.
    \end{equation*}
    Außerdem gilt $C(\pi^*) \geq\displaystyle\max_{j \in J} p_j$.

    Sei $\pi$ der Schedule, der \textsc{LeastLoaded} berechnet. Sei $i$ die Maschine mit größter Last. Sei $j$ der Job, der zuletzt zu $i$ hinzugefügt wurde. Zu diesem Zeitpunkt war Maschine $i$ diejenige mit der geringsten Last. Ihre Last war höchstens $\frac{1}{m}\sum_{k = 1}^{j-1} p_k$. Es gilt 
    \begin{equation*}
        C(\pi) = L_i(\pi) \leq \frac{1}{m}\underset{\leq\sum_{k \in J \setminus \set{j}} p_k}{\underbrace{\sum_{k = 1}^{j-1} p_k}} + p_j \leq \underset{\leq C(\pi^*)}{\underbrace{\frac{1}{m} \sum_{k \in J} p_k}} + \left( 1 - \frac{1}{m} \right) \underset{\leq C(\pi^*)}{\underbrace{p_j}} \leq \left( 2 - \frac{1}{m} \right)C(\pi^*)
    \end{equation*}
    und die Abschätzung ist gezeigt.
\end{proof}

\begin{rmk}
    $2-\frac{1}{m}$ ist eine untere Schranke für \textsc{LeastLoaded}.
    \begin{proof}
        Sei $n = m(m-1) + 1$ und $p_1 = \dots = p_{m(m-1)} = 1$ und $p_{m(m-1)+1} = m$. 

        \textsc{LeastLoaded} führt auf jeder Maschine $m-1$ Jobs der Größe $1$ durch. Auf einer Maschine wird zusätzlich ein Job der Größe $m$ ausgeführt. Somit ist der Makespan $m-1+m = 2m-1$.

        Eine optimale Belegung führt auf $m-1$ Maschinen je $m$ Jobs der Größe $1$ aus und auf der letzten Maschine einen Job der Größe $m$, wir haben also einen Makespan von $m$, und der Faktor stimmt.
    \end{proof}
\end{rmk}

\begin{algorithm}
    \caption{\textsc{LongestProcessingTime} (LPT)}\index{LongestProcessingTime}\index{LPT}
    \begin{algorithmic}[1]
        \Input{Siehe~\Cref{prob:SchedulingIdenticalMachines}}
        \Output{Siehe~\Cref{prob:SchedulingIdenticalMachines}}
        \State{Sortiere Jobs so, dass $p_1 \geq p_2 \geq \dots \geq p_n$.}
        \State{Führe \textsc{LeastLoaded} auf den sortierten Jobs aus.}
    \end{algorithmic}
\end{algorithm}
\begin{thm}{LPT}{LPT}
    \textsc{LongestProcessingTime} ist ein $\frac{4}{3}$-Approximationsalgorithmus.
\end{thm}
\begin{proof}
    Angenommen, es gäbe eine Instanz, sodass LPT einen Makespan von mehr als $\frac{4}{3} \cdot \OPT$ erreicht. Betrachte unter diesen Instanzen eine mit geringster Anzahl Jobs.
    
    Sei $i$ eine Maschine mit höchster Last in LPT-Schedule $\pi$. Sei $j$ der letzte Job, der Maschine $i$ zugewiesen wird. Es ist $j = n$, da sonst $p_1, \ldots, p_j$ ein Gegenbeispiel mit weniger Jobs wäre. Es gilt 
    \begin{equation*}
        C(\pi) = L_i(\pi) \leq \frac{1}{m}\sum_{k=1}^{n-1} p_k + p_n \leq \OPT + p_n.
    \end{equation*}
    Aus $C(\pi) > \frac{4}{3}\OPT$ folgt $p_n > \frac{1}{3}\OPT$. Da $p_1 \geq \dots \geq p_n$ folgt $p_j > \frac{1}{3}\OPT$ für alle $j \in J$. Also erhält im optimalen Schedule jede Maschine höchstens zwei Jobs. 
    
    Sei $k$ die Anzahl Jobs mit Größe $\frac{2}{3}\OPT$. Im optimalen Schedule sind $k$ Jobs alleine auf einer Maschine und $n-k$ Jobs höchstens zu zweit. 

    \begin{description}
        \item[Fall 1: $k = n$.] Also ist $n \leq m$, Job $n$ wird also einer leeren Maschine zugewiesen. 
        \item[Fall 2: $k < n$.] Zum Zeitpunkt bevor Job $n$ hinzugefügt wird, hat LPT mindestens eine Maschine mit höchstens einem Job und dieser hat Größe höchstens $\frac{2}{3}\OPT$. Wegen $p_n \leq \frac{2}{3}\OPT$ ist die Last dieser Maschine im Anschluss höchstens $\frac{4}{3}\OPT$.
    \end{description}
    Beides ist ein Widerspruch zur Annahme. $\lightning$
\end{proof}

\begin{rmk}
    $\frac{4}{3}$ ist eine untere Schranke für \textsc{LongestProcessingTime}.
    \begin{proof}
        Übungsaufgabe.
    \end{proof}
    Wir werden später sehen, dass für alle $\varepsilon > 0$ ein $(1+\varepsilon)$-Approximationsalgorithmus existiert.
\end{rmk}


\newLecture{21.10.2025}

\section{Rucksack-Problem}

\newLecture{22.10.2025}

\end{document}